\subsection{Truncations}\label{sec:truncation}



\begin{definition}\label{def:n-step}
A cubeset $X$ is of \emph{degree $n$} (or \emph{$n$-step}) for $n\geq 0$ if corner completion strictly above dimension $n$ is unique.
This means that for every $k > n$ restriction along the $k$-corner
\[
    \hom({\bbox}^k,X)\to\hom({\bbcor}^k,X)
\]
is an isomorphism (bijective).
A cubeset $X$ is of \emph{degree $-1$} if it is either empty or the point,
and it is of \emph{degree $-2$} if it is the point.
We denote the full subcategory of cubesets of degree $n$ by $\cSet_n$.
\end{definition}

\begin{remark}
If we let $\Lambda_n = \{{\bbcor}^k\to{\bbox}^k : k > n\}$, then the cubesets of degree $n$ are precisely the $\Lambda_n$-local objects.
In particular, a $\Lambda_n$-local object $X$ cannot distinguish between a $k$-corner and a $k$-cube.
\end{remark}

\begin{remark}\label{rem:corner-no-topology}
Note that, however, $\Lambda_n$ does not induce a Lawvere-Tierney topology on $\cSet$, or equivalently a Grothendieck topology on $\bbox$.
If it were, then to any morphism $f\colon{\bbox}^n\to{\bbox}^m$, a cover of ${\bbox}^m$ would allow a pullback along $f$ to a cover of ${\bbox}^n$.
But if $1^m\colon\ast\to{\bbox}^m$ is the point $1^m$ in $\{0,1\}^m$, the corner in ${\bbox}^m$ does not factor through this point (which it would have to, since the point only is covered by the point).
\end{remark}

\begin{proposition}\label{prop:truncations-char}
Let $X$ be a cubeset. Then the following assertions are equivalent.
\begin{enumerate}[(a)]
	\item The cubeset $X$ is $n$-step, i.e. the pullback map
	\[
        \hom({\bbox}^k,X)\to\hom({\bbcor}^k,X)
    \]
    is bijective for all $k>n$.
	\item For every $f\colon Y\to Z\in l(r(\Lambda_n))$
    (i.e., a trivial cofibration/retract of a relative cell complex w.r.t. corner inclusions of dimension $\geq n+1$),
	one has that
    \[
        f^\ast\colon\hom(Z,X)\to \hom(Y,X)
    \]
    is a bijection.
    \item It is \emph{internally $n$-step}, meaning that
	\[
        \ihom({\bbox}^k,X)\to\ihom({\bbcor}^k,X)
    \]
	is an isomorphism of cubesets for all $k>n$.
    \item For every $f\colon Y\to Z\in l(r(\Lambda_n))$
	(i.e. a trivial cofibration/retract of a relative cell complex w.r.t. corner inclusions of dimension $\geq n+1$),
    one has that
	\[
    	\ihom(Z, X)\to \ihom(Y,X)
	\]
	is an isomorphism.
    \item The map
	\[
        \ihom({\bbox}^{n+1},X)\to\ihom({\bbcor}^{n+1},X)
    \]
    is an isomorphism.
\end{enumerate}
\end{proposition}

\begin{proof}
	A morphism of cubesets is an isomorphism if and only if it is a bijection pointwise.
	This means that a cubeset $X$ is internally $n$-step precisely if
    \[
        \hom({\bbox}^k\times{\bbox}^j,X)\to\hom({\bbcor}^k\times{\bbox}^j,X)
    \]
	is a bijection for all $j\in \N_0$ and all $k>n$.
	Thus, the internal variants are stronger than the non-internal variants (being the internal ones at $j=0$).

	$(a)\implies(b)$ and $(c)\implies(d)$: It suffices to see that the preimage of the class of isomorphisms under a cocontinuous functor is saturated.
	For this note that the pushout along an isomorphism is again an isomorphism,
	as well as transfinite compositions of isomorphisms and retracts.

    $(a)\implies(c)$:
	We follow the route of Remark~\ref{rem:internal-nilfib}.
	In Corollary~\ref{cor:internal-nilfib} we have shown that the roof $\bbcor^{k}\times{\bbox}^j\to{\bbox}^{k+j}$
	is a relative cell complex built from corner inclusions ${\bbcor}^i\to{\bbox}^i$ for $k\le i$
	(one examines the proof carefully to see that there is no other $i$ involved, see Remark~\ref{prop:internal-nilfib}).
	This implies that the roofs ${\bbcor}^k\times{\bbox}^j\to{\bbox}^{k+j}$
	are in the saturated closure of all ${\bbcor}^i\to{\bbox}^i$ for $k\le i$.
	Thus we know that having unique corner completion implies unique internal corner completion.

    $(e)\implies(a)$:
    We need to show that the natural map
    \[
        \hom({\bbox}^{n+1+j},X)\to\hom({\bbcor}^{n+1+j},X)
    \]
    is bijective for all $j\in\N_0$.
    We have a factorization
    \begin{center}
    \begin{tikzcd}
	{\hom({\bbox}^{n+1+j},X)} && {\hom({\bbcor}^{n+1}\times{\bbox}^j,X)} \\
	& {\hom({\bbcor}^{n+1+j},X)}
	\arrow[from=1-1, to=1-3]
	\arrow[from=1-1, to=2-2]
	\arrow[from=2-2, to=1-3]
    \end{tikzcd}
    \end{center}
    where the horizontal map is bijective for all $j\in\N_0$ by assumption on the internal Hom.
    For $j=0$, the left map is bijective by assumption (since then the right map degenerates).
    In the induction step, suppose both the left and right map are bijections.
    Now the map ${\bbcor}^{n+1+j}\times{\bbox}^1\to{\bbcor}^{n+1+j+1}$ is a pushout
    along the ($n+1+j$)-dimensional corner inclusion ${\bbcor}^{n+1+j}\to{\bbox}^{n+1+j}$
    (see the proof of Proposition~\ref{prop:internal-nilfib}),
    so pullback along the former induces a bijection on $\hom$-sets.
    Moreover, by induction hypothesis, the map ${\bbcor}^{n+1}\times{\bbox}^j\to{\bbcor}^{n+1+j}$
    induces a bijection on the level of $\hom(-,X)$.
    Thus we can use the factorization
    \begin{center}
    \begin{tikzcd}
	{{\bbcor}^{n+1}\times{\bbox}^{j+1}} & {{\bbcor}^{n+1+j}\times{\bbox}^1} & {{\bbcor}^{n+1+j+1}}
	\arrow[from=1-1, to=1-2]
	\arrow[from=1-2, to=1-3]
    \end{tikzcd}
    \end{center}
    to see that the composite induces a bijection on $\hom(-,X)$.
    In turn this means that both $\hom({\bbcor}^{n+1+j},X)\to\hom({\bbcor}^{n+1}\times{\bbox}^j,X)$
    and hence $\hom({\bbox}^{n+1+j},X)\to\hom({\bbcor}^{n+1+j},X)$ are bijections, completing the induction.
\end{proof}

\begin{remark}
	Note that imposing the bijection for a single $k$ externally fails:
	That
	\[\hom({\bbox}^k,X)\to \hom({\bbcor}^k,X)\]
	is bijective is not equivalent to the map
    \[\ihom({\bbox}^k,X)\to \ihom({\bbcor}^k,X)\]
    being an isomorphism,
    since the latter considers all dimensions.
	As an explicit counterexample take some cubeset with $X(0)=X(1)=\ast$ but $X(2)\ne \ast$ and consider
	$k=1$.
\end{remark}


\begin{proposition}\label{prop:step-left-adjoint}
The inclusion of $n$-step cubesets into all cubesets admits a left adjoint which we call the canonical truncation $\tau_n\colon\cSet\to\cSet_n$.
\end{proposition}

\begin{proof}
This follows from \cite[Theorem 1.39]{Adamek1994}.
Alternatively, one can use the $\infty$-categorical version from \cite[Proposition 5.5.4.2]{Lurie2009}.
\end{proof}

First, let us remark that a certain fibrant replacement does not change the truncation.

\begin{lemma}
	Let $X\to X^{\mathrm{fib},n}$ be a fibrant replacement of $X$ starting in dimension $n+1$,
	i.e., a fibrant replacement w.r.t. $\Lambda_n$.
	Then
	\[
	    \tau_{n}X=\tau_n (X^{\mathrm{fib},n}).
	\]
\end{lemma}

\begin{proof}
The map $X\to X^{\mathrm{fib},n}$ is a trivial cofibration w.r.t. $\Lambda_n$,
i.e., $(X\to X^{\mathrm{fib},n})\in l(r(\Lambda_n))$.
By Proposition~\ref{prop:truncations-char}~(b), for every $n$-step cubeset $Y$, there is a bijection
\[
    \hom(X^{\mathrm{fib},n},Y) \to \hom(X,Y)
\]
natural in $X$ and $Y$.
Applying the truncation to both sides, the result follows from the Yoneda lemma.
\end{proof}

Thus, a good first step to truncate is to build a certain fibrant replacement.
Note that fibrant replacement a priori is not a left adjoint due to the impossibility of choosing a canonical corner completion.
Testing against objects with unique corner completion however makes this lift unique and thus we obtain the left adjointness.

\begin{remark}
Note that truncation does not preserve finite limits in general.
If it would do so, then $\cSet_n\to\cSet$ would be a geometric morphism
with the inclusion corresponding to a subtopos on $\cSet$.
Hence there exists a Grothendieck topology on $\bbox$ such that the $n$-step cubesets are the sheaves for this topology:
the sieves are precisely the injections $f\colon S\to{\bbox}^k$ for which pullback $f^\ast\colon\hom({\bbox}^k,X)\to\hom(S,X)$
is an isomorphism for each $n$-step cubeset $X$ (and $S$ corresponds to a subfunctor of the representable ${\bbox}^k$).
In particular, the corner inclusion ${\bbcor}^{n+1}\to{\bbox}^{n+1}$ is a cover for this topology.
But then Remark~\ref{rem:corner-no-topology} yields a contradiction to stability of covers under pullback.
\end{remark}

However, the truncation does preserve finite products.

\begin{proposition}\label{prop:n-trunc-ihom-prod}
If $X$ and $Y$ are cubesets and $Y$ is $n$-step, then $\ihom(X,Y)$ is $n$-step.
Moreover, the $n$-truncation preserves finite products.
\end{proposition}

\begin{proof}
First we show that $\ihom(X,Y)$ is $n$-step if $Y$ is so.
We need to see that pullback along the inclusion ${\bbcor}^k\to{\bbox}^k$ induces a bijection
\[
    \hom({\bbox}^k\times X,Y) = \hom({\bbox}^k,\ihom(X,Y)) \to \hom({\bbcor}^k,\ihom(X,Y)) = \hom({\bbcor}^k\times X,Y)
\]
for all $k>n$.
For $X = {\bbox}^m$ it follows from Corollary~\ref{cor:internal-nilfib}
and Remark~\ref{prop:internal-nilfib},
see also the proof of $(a)\implies(c)$ of Proposition~\ref{prop:truncations-char}.
But then it follows for general $X$ by writing $X$ as the colimit of representables,
using that colimits are stable under base change in $\cSet$.

For the second claim, one argues as in the last part of the proof of
Proposition~\ref{prop:concrete-ihom-prod}.
\end{proof}

The category of $n$-step cubesets additionally has the following closure properties.

\begin{proposition}\label{prop:n-trunc-closure}
The category of $n$-step cubesets $\cSet_n$ is closed under limits, filtered colimits and coproducts in $\cSet$.
\end{proposition}

\begin{proof}
	By continuity of the Hom-functor, $\cSet_n$ is closed under limits.
Closure under filtered colimits follows from the fact that ${\bbcor}^k$ is compact for all $k\in\N$.
Closure under coproducts follows from the fact that ${\bbcor}^k$ is connected for all $k\in\N$.
\end{proof}



\subsubsection{Computing the Truncation of Concrete Fibrant Cubesets}\label{ssec:truncation-compute}

The localisation $\tau_{n}X$ may always be computed naively by some iterated construction, see \cite[Construction 1.37]{Adamek1994}:
\begin{itemize}
	\item Glue to $X$ a $k$-cube for every $k$-corner which does not admit a filler in $X$ (for all $k> n$).
	\item Coequalize two $k$-cubes if they agree on their $k$-corner (for all $k>n$).
	\item Iterate this process until it stabilizes (at limit ordinals take the colimit).
\end{itemize}
In our case, the morphisms at which we localize have compact domain and codomain
and thus the process stabilizes after $\aleph_0$-many steps.
For general cubespaces, all these steps and in particular such an iterated construction seem unavoidable,
as the following examples show.
\begin{example}
\begin{enumerate}[(i)]
	\item For an example where one needs iteration of the identification step consider the double Kebab bread
	 \begin{center}
\begin{tikzcd}[ampersand replacement=\&, row sep=tiny, column sep=tiny]
	\&\&\&\&\&\& \bullet \& \\
	\& \bullet \& \bullet \&\&\&\&\& \bullet \\
	\\
	\bullet \&\&\& \bullet \\
	\&\&\&\& \bullet \\
	\&\&\&\&\&\& \bullet \\
	\& \bullet \&\&\&\&\& \bullet \\
	\&\&\&\& \bullet
	\arrow[dotted, no head, from=2-2, to=1-7]
	\arrow[dashed, no head, from=2-2, to=4-4]
	\arrow[dotted, no head, from=2-3, to=2-8]
	\arrow[dashed, no head, from=2-3, to=4-4]
	\arrow[dashed, no head, from=4-1, to=2-2]
	\arrow[dashed, no head, from=4-1, to=2-3]
	\arrow[dashed, no head, from=5-5, to=6-7]
	\arrow[dashed, no head, from=5-5, to=7-7]
	\arrow[dotted, no head, from=6-7, to=1-7]
	\arrow[no head, from=7-2, to=1-7]
	\arrow[no head, from=7-2, to=2-2]
	\arrow[no head, from=7-2, to=2-3]
	\arrow[no head, from=7-2, to=2-8]
	\arrow[no head, from=7-2, to=4-1]
	\arrow[no head, from=7-2, to=4-4]
	\arrow[no head, from=7-2, to=5-5]
	\arrow[no head, from=7-2, to=6-7]
	\arrow[no head, from=7-2, to=7-7]
	\arrow[no head, from=7-2, to=8-5]
	\arrow[dotted, no head, from=7-7, to=2-8]
	\arrow[dashed, no head, from=8-5, to=6-7]
	\arrow[dashed, no head, from=8-5, to=7-7]
\end{tikzcd}
\end{center}
	\item For the necessity of the insertion step simply take a non-fibrant cubeset with many corners that do not yet exist, but assure that if they exist they are unique (e.g., take the $k$-skeleton of a nontrivial concrete fibrant cubeset).
	\item To see that in each identification step one needs to take the induced equivalence relation consider
	\begin{center}
\begin{tikzcd}[ampersand replacement=\&, row sep= tiny, column sep=tiny]
	\&\&\&\&\& b \& \\
	\&\&\& d \&\&\& c \\
	\\
	\& f \&\& a \&\& e \\
	\\
	d \&\&\& g \\
	\& c
	\arrow[dashed, no head, from=1-6, to=4-6]
	\arrow[dashed, no head, from=2-4, to=1-6]
	\arrow[dashed, no head, from=2-7, to=2-4]
	\arrow[no head, from=4-2, to=4-4]
	\arrow[dashed, no head, from=4-2, to=6-1]
	\arrow[dashed, no head, from=4-2, to=7-2]
	\arrow[no head, from=4-4, to=1-6]
	\arrow[no head, from=4-4, to=2-4]
	\arrow[no head, from=4-4, to=2-7]
	\arrow[no head, from=4-4, to=4-6]
	\arrow[no head, from=4-4, to=6-1]
	\arrow[no head, from=4-4, to=6-4]
	\arrow[no head, from=4-4, to=7-2]
	\arrow[dashed, no head, from=4-6, to=2-7]
	\arrow[dashed, no head, from=6-1, to=6-4]
	\arrow[dashed, no head, from=6-4, to=7-2]
\end{tikzcd},

		\end{center}
	where $b\sim c$ and $d\sim c$ but there is no cube witnessing $b\sim d$.
	\item To see that \emph{first identify all necessary cubes and afterward glue all that you need} does not suffice, consider
	 \begin{center}
\begin{tikzcd}[ampersand replacement=\&, row sep=tiny, column sep=tiny]
	\&\&\& {h_1} \& {h_2} \\
	\& e \&\&\& {h_3} \\
	b \&\& g \\
	\& c \&\& f \\
	a \&\& d
	\arrow[no head, from=2-2, to=1-4]
	\arrow[no head, from=2-2, to=1-5]
	\arrow[no head, from=3-1, to=2-2]
	\arrow[no head, from=3-1, to=3-3]
	\arrow[no head, from=3-3, to=1-4]
	\arrow[no head, from=3-3, to=2-5]
	\arrow[no head, from=4-2, to=2-2]
	\arrow[no head, from=4-2, to=4-4]
	\arrow[no head, from=4-4, to=1-5]
	\arrow[no head, from=4-4, to=2-5]
	\arrow[no head, from=5-1, to=3-1]
	\arrow[no head, from=5-1, to=4-2]
	\arrow[no head, from=5-1, to=5-3]
	\arrow[no head, from=5-3, to=3-3]
	\arrow[no head, from=5-3, to=4-4]
\end{tikzcd}.
\end{center}
Here, the three-dimensional corner has no filler.
If it had one, the new vertex would simultaneously be 2-equivalent to $h_1$, $h_2$ and $h_3$ and thus inserting the filler makes the $h_i$ equivalent.
 \end{enumerate}
\end{example}

Also note that concreteness is not preserved by the truncation.
However, together with fibrancy and concreteness, we may significantly simplify the computation of the truncation.

\begin{definition}\label{def:univ-replacement}
For $x,y\in X(0)$ define $x\sim_n y$ iff there exist $c,c'\in X(n+1)$ such that
\[
    c\rvert_{{\bbcor}^{n+1}} = c'\rvert_{{\bbcor}^{n+1}} \quad\text{and}\quad c(1) = x,\; c'(1) = y.
\]
\end{definition}

We can characterize this relation further as follows:

\begin{proposition}\label{prop:univ-replacement}
Let $X$ be a concrete fibrant cubeset, $n\in\N_0$ and $x,y\in X(0)$.
Then the following assertions are equivalent.
\begin{enumerate}[(a)]
    \item $x\sim_n y$,
    \item The configuration $c\rvert_{{\bbcor}^{n+1}} \equiv x$, $c(1) = y$ is in $X(n+1)$.
    \item For every $k\leq n+1$ and every cube $c\in X(k)$ with value $x$ at some vertex,
            the same cube with value $y$ at that vertex is in $X(k)$.
\end{enumerate}
In particular, the relation $\sim_n$ is an equivalence relation on $X(0)$.
\end{proposition}

This can e.g. be found in \cite{Gutman2020}
and follows mainly from the fact that every concrete fibrant cubeset has the \emph{gluing property}.
The property in (c) is also referred to as \emph{universal replacement}.
Using this we see that $x\sim_n y$ is an equivalence relation on $X(0)$, inducing a factor $X(0)/\sim_n$.
This yields a new cubeset $X/\sim_n$ whose value $X/\sim_n(m)$ is given by the image of $X(m)$ under the projection
\[
    X(0)^{2^m}\to (X(0)/\sim_n)^{2^m}.
\]
Explicitly this means that two $m$-cubes in $X$ are identified if all their points are equivalent in $X(0)$.
This cubeset thus is a factor of $X$ and allows for an easy description of $\tau_n X$, as the following proposition shows.

\begin{proposition}\label{prop:trunc_concfib}
If $X$ is a concrete fibrant cubeset, then the truncation $\tau_n X$ is given by $X/\sim_n$ and the unit
$\eta_X\colon X\to \tau_n X$ is a surjective fibration.
In particular, $\tau_n X$ is concrete and fibrant.
\end{proposition}

\begin{proof}
Denote $\eta\colon X\to\tau_n X$ the unit of the $n$-step adjunction and $\pi\colon X\to X/\sim_n$ the factor map.
Let us begin with establishing the existence of a morphism $f\colon X/\sim_n\to\tau_n X$ such that $f\circ\pi = \eta$, i.e., $\eta$ factors through $\pi$.
To do so we have to show that if $a$ and $b$ are two $k$-cubes in $X$ which are $n$-equivalent, i.e., $a(\omega)\sim_n b(\omega)$ for all $\omega\in\{0,1\}^k$,
then $\eta(a) = \eta(b)$.

If $k\leq n+1$ take a surjection $p\colon{\bbox}^{n+1}\to{\bbox}^k$.
Then $a' = p^\ast(a)\sim_n p^\ast(b) = b'$ in $X(n+1)$.
This allows us to first use universal replacement, and afterwards unique corner completion in $\tau_n X$.
Using universal replacement, the configuration
\[
    c_0\colon\{0,1\}^{n+1}\to X(0),\quad\omega\mapsto\begin{cases}
        b'(0),\quad &\omega = 0, \\
        a'(\omega),\quad &\omega\neq 0,
    \end{cases}
\]
is an element of $X(n+1)$ since $a'(0)\sim_n b'(0)$.
But after an appropriate rotation, $a'\rvert_{{\bbcor}^{n+1}} = c_0\rvert_{{\bbcor}^{n+1}}$
and so $\eta(a')\rvert_{{\bbcor}^{n+1}} = \eta(c_0)\rvert_{{\bbcor}^{n+1}}$.
Using unique corner completion in $\tau_n X$, we obtain that $\eta(a') = \eta(c_0)$.
Iterating this process, exchanging $a'(\omega)$ for $b'(\omega)$ one after another yields a sequence $a',c_0,c_1,\ldots,b'$ of ($n+1$)-cubes in $X$
with $\eta(a') = \eta(c_0) = \eta(c_1) = \ldots = \eta(b')$.
Taking the corresponding face of $a'$ and $b'$ we obtain that $\eta(a) = \eta(b)$.

For $k>n+1$, we induct on $k$: the induction start for $k = n+1$ has already been established.
If $k \geq n+2$ and $a\sim_n b$ for $a,b\in X(k)$,
then for every inert $\epsilon_i\colon{\bbox}^{k-1}\to{\bbox}^{k}$ we have that $\epsilon_i^\ast(a) \sim_n \epsilon_i^\ast(b)$ in $X(k-1)$.
By induction hypothesis, $\eta(\epsilon_i^\ast(a)) = \eta(\epsilon_i^\ast(b))$ which in turn yields $\eta(a)\rvert_{{\bbcor}^{k}} = \eta(b)\rvert_{{\bbcor}^{k}}$.
Using unique corner completion in $\tau_n X$ we obtain that $\eta(a) = \eta(b)$.
This gives the desired factorization $f\circ\pi = \eta$.

Next, let us prove that the cubeset $X/\sim_n$ is $n$-step.
For this we first show that each corner in $X/\sim_n$ of dimension $k > n$ has at most one completion.
Since $X/\sim_n$ is concrete, it suffices to consider the case $k = n+1$ (otherwise, take a diagonal ${\bbox}^{n+1}\to{\bbox}^k$ mapping the corner into the corner and hitting the point $1^k$).
So take an ($n+1$)-corner $\lambda$ in $X/\sim_n$ with completions $a$ and $b$.
Take $a',b'\in X(n+1)$ with $\pi(a') = a$ and $\pi(b') = b$.
Then we obtain that $\pi(a')\rvert_{{\bbcor}^{n+1}} = \lambda = \pi(b')\rvert_{{\bbcor}^{n+1}}$.
This means that $a'(\omega) \sim_n b'(\omega)$ for all $\omega\neq 1^{n+1}$.
Using universal replacement $2^{n+1}-1$ many times, the configuration
\[
    \{0,1\}^{n+1}\to X(0),\quad\omega\mapsto\begin{cases}
        a'(\omega),\quad &\omega\neq 1^{n+1}, \\
        b'(1^{n+1}),\quad &\omega = 1^{n+1},
    \end{cases}
\]
is in $X(n+1)$.
But this means that $a'(1^{n+1}) \sim_n b'(1^{n+1})$ and so $a(1^{n+1}) = b(1^{n+1})$.
By concreteness, $a = b$.

Next we prove that the map $\pi\colon X\to X/\sim_n$ is a fibration.
\begin{center}
\begin{tikzcd}
	{{\bbcor}^k} & X \\
	{{\bbox}^k} & {X/\sim_n}
	\arrow[from=1-1, to=1-2]
	\arrow[hook', from=1-1, to=2-1]
	\arrow["\pi", from=1-2, to=2-2]
	\arrow[dashed, from=2-1, to=1-2]
	\arrow[from=2-1, to=2-2]
\end{tikzcd}
\end{center}
For $k\leq n+1$ the lift exists by universal replacement: first choose a cube $c$ in $X$ that maps to ${\bbox}^k\to X/\sim_n$
(this is possible by surjectivity of $\pi$) and then replace all points of $c$ in the corner with the points of the given corner ${\bbcor}^k\to X$,
which are equal in $X/\sim_n$.
For $k\geq n+1$ again complete the corner arbitrarily in $X$.
Uniqueness of corner completion in $X/\sim_n$ asserts that this completion pushes down to the cube we started with in $X/\sim_n$.

In particular, $X/\sim_n$ is a fibrant cubeset:
\begin{center}
\begin{tikzcd}
	\ast & X \\
	{{\bbcor}^k} & {X/\sim_n} \\
	{{\bbox}^k} & \ast
	\arrow[from=1-1, to=1-2]
	\arrow[hook', from=1-1, to=2-1]
	\arrow["\pi", from=1-2, to=2-2]
	\arrow[dashed, from=2-1, to=1-2]
	\arrow[from=2-1, to=2-2]
	\arrow[hook', from=2-1, to=3-1]
	\arrow["p", from=2-2, to=3-2]
	\arrow[dashed, from=3-1, to=1-2]
	\arrow[from=3-1, to=3-2]
\end{tikzcd}
\end{center}
since $\pi$ is an epic fibration one can lift the $k$-corner in $\tau_nX$ to $X$ (inductively completing lower dimensional corners) and then use that $X$ is fibrant to complete to a $k$-cube in $X$ which pushes down to the desired completion.

Taking the previous steps together, we obtain that $X/\sim_n$ is a concrete fibrant $n$-step cubeset.
By adjunction, there exists a unique map $g$ making the diagram
\begin{center}
\begin{tikzcd}
	X & {\tau_nX} \\
	& {X/\sim_n}
	\arrow["\eta", from=1-1, to=1-2]
	\arrow["\pi"', from=1-1, to=2-2]
	\arrow["g", dashed, from=1-2, to=2-2]
\end{tikzcd}
\end{center}
commutative.
This means that $f\circ g = \id$ since $\eta$ is the unit of the adjunction, and $g\circ f = \id$ since $\pi$ is surjective.
So we have shown that $\tau_n X = X/\sim_n$.
\end{proof}

\begin{remark}
Note that in particular this implies that the equivalence relation of the $n$-truncation on $n+1$-cubes is simply generated by identifying two $n+1$-cubes if they have the same corner,
i.e., it identifies with the coequalizer of both arrows $\hom(\bbox^{n+1}\sqcup_{\bbcor^{n+1}}\bbox^{n+1},X)\to \hom(\bbox^{n+1},X)=X(n+1)$.
\end{remark}

\begin{remark}
The previous proposition shows that the filtered category $(\cSet_n)_n$ restricts to a filtration on concrete fibrant cubesets.
\begin{center}
\begin{tikzcd}
	\cSet & \cdots & {\cSet_2} & {\cSet_1} & {\cSet_0} \\
	\ccSetnil & \cdots & {\cSet^\mathrm{conc}_{\mathrm{fib},2}} & {\cSet^\mathrm{conc}_{\mathrm{fib},1}} & {\cSet^\mathrm{conc}_{\mathrm{fib},0}}
	\arrow[shift left, from=1-1, to=1-2]
	\arrow[shift left, hook, from=1-2, to=1-1]
	\arrow["{\tau_2}", shift left, from=1-2, to=1-3]
	\arrow[shift left, hook, from=1-3, to=1-2]
	\arrow["{\tau_1}", shift left, from=1-3, to=1-4]
	\arrow[shift left, hook, from=1-4, to=1-3]
	\arrow["{\tau_0}", shift left, from=1-4, to=1-5]
	\arrow[shift left, hook, from=1-5, to=1-4]
	\arrow[hook, from=2-1, to=1-1]
	\arrow[shift left, from=2-1, to=2-2]
	\arrow[shift left, hook, from=2-2, to=2-1]
	\arrow["{\tau_2}", shift left, from=2-2, to=2-3]
	\arrow[hook, from=2-3, to=1-3]
	\arrow[shift left, hook, from=2-3, to=2-2]
	\arrow["{\tau_1}", shift left, from=2-3, to=2-4]
	\arrow[hook, from=2-4, to=1-4]
	\arrow[shift left, hook, from=2-4, to=2-3]
	\arrow["{\tau_0}", shift left, from=2-4, to=2-5]
	\arrow[hook, from=2-5, to=1-5]
	\arrow[shift left, hook, from=2-5, to=2-4]
\end{tikzcd}
\end{center}\end{remark}

\begin{example}\label{ex:truncation_HK}
The $n$-truncation of a Host-Kra cubegroup associated to a filtered group $G_\bullet$ is given by passing to the filtration quotiented by $G_{n+1}$,
i.e. taking the filtered group defined by $G_0/G_{n+1}\ge G_1/G_{n+1}\ge\cdots$.
Thus, the length of the filtration is the step of the cubeset.

In particular, the $n$-truncation of $\cD_k(A)$ is $0$ for $n< k$,
and is $\cD_k(A)$ for all $n\ge k$.
Hence, $\cD_k(A)$ is $k$-truncated.
\end{example}

\begin{remark}
There are several different and more categorical constructions that yield the truncation in the concrete fibrant setting, but need iteration and the gluing step in the general case.
As an example of such a construction, consider the two canonical maps
\[
    \ihom({\bbox}^{n+1}\sqcup_{{\bbcor}^{n+1}}{\bbox}^{n+1},X)\times{\bbox}^{n+1}\to X,
\]
which are obtained from currying the two canonical maps
\[
    \ihom({\bbox}^{n+1}\sqcup_{{\bbcor}^{n+1}}{\bbox}^{n+1},X)\to \ihom({\bbox}^{n+1},X).
\]
Taking the coequalizer
\begin{center}
\begin{tikzcd}[ampersand replacement=\&]
	{\ihom(\bbox^{n+1}\sqcup_{\bbcor^{n+1}}\bbox^{n+1},X)\times {\bbox}^{n+1}} \&\& X \&\& {Y}
	\arrow["{\phi_2}"', shift right=1, from=1-1, to=1-3]
	\arrow["{\phi_1}", shift left=1, from=1-1, to=1-3]
	\arrow[two heads, from=1-3, to=1-5]
\end{tikzcd}
\end{center}
gives a quotient $Y$.
If $Y$ is $n$-step, then it agrees with the $n$-truncation of $X$.
\end{remark}

\subsubsection{Connectivity}\label{ssec:connectivity}

We further investigate $0$-step cubesets and talk about connectivity.
Recall the adjunction $i_!\dashv i^\ast$ from Lemma~\ref{lem:point-cset}.

\begin{proposition}\label{prop:0-trunc-computation}
The adjunction $i_!\dashv i^\ast$ induces an equivalence of categories $\set\simeq\cSet_0$.
In particular, a cubeset $X$ is $0$-step if and only if the counit $i_!i^\ast X\to X$ is an isomorphism.
Under this equivalence, the $0$-truncation $\tau_0\colon\cSet\to\set$ is given by
\[
    X\mapsto\pi_0(X) \coloneq (\tau_0X)(0) = \varinjlim_{n\in{\bbox}^\op}X(n).
\]
In particular, a cubeset $X$ is $0$-step if and only if the unit $X\to i_!\pi_0(X)$ is an isomorphism.
Moreover, the diagram
\begin{center}
\begin{tikzcd}
	{X(1)} & {X(0)} & {\pi_0(X)}
	\arrow["{s^\ast}", shift left, from=1-1, to=1-2]
	\arrow["{t^\ast}"', shift right, from=1-1, to=1-2]
	\arrow["\eta", from=1-2, to=1-3]
\end{tikzcd}
\end{center}
is a coequalizer diagram, where $s,t$ are the two maps ${\bbox}^0\to{\bbox}^1$.
\end{proposition}

\begin{proof}
Since the left adjoint is fully faithful, the adjunction induces an equivalence of categories
between $\set$ and the full subcategory of $\cSet$ spanned by objects $X$ such that the
counit $i_!i^\ast X\to X$ is an isomorphism.
Thus it suffices to show that the condition on the counit is equivalent to $X$ being $0$-step.
If $X\simeq i_!i^\ast X$, then $X$ is a constant presheaf which clearly has unique corner completion
in all dimensions and so is $0$-step.
In the other direction, assume that $X$ is $0$-step.
Recall that we can write
\[
    {\bbcor}^k = \varinjlim_{j\in\cJ}{\bbox}^j
\]
where $\cJ$ runs over all inert inclusions ${\bbox}^j\to{\bbox}^k$ for $j<k$.
Thus we have isomorphisms
\[
    X(k)\simeq\varprojlim_{j\in\cJ}X(j)
\]
for all $k\geq 1$.
By induction, we infer that $X(k)\simeq X(0)$ for all $k\geq 1$.
Thus, $X\simeq i_!i^\ast X$.

For the second claim note that if $X$ is a cubeset and $Y$ is $0$-step
we have by the previous considerations that
\[
    \hom(X,Y) = \hom(X,i_!i^\ast Y) = \hom\left(\varinjlim_{n\in{\bbox}^\op} X(n),Y(0)\right)
\]
where the second equality follows from Remark~\ref{rem:left-adjoint-i_!}.
Here, the cubeset $X$ is interpreted as a diagram $X\colon{\bbox}^\op\to\set$
whose colimit in $\set$ exists.
This means that $(\tau_0X)(0) = \varinjlim_{n\in{\bbox}^\op} X(n)$.

For the last assertion, given a map $f\colon X(0)\to S$ with $f\circ s^\ast = f\circ t^\ast$
we need to show that there is a map $g_0\colon\pi_0(X)\to S$ such that $f = g_0\circ\eta$.
By the adjunction $\pi_0(X)\dashv i_!$, this is equivalent to showing that there exists a map
$g\colon X\to i_!S$ such that $g_0 = f\colon X(0)\to i_!S(0)$.
Since $i_!S$ is concrete it suffices to assume that $X$ is concrete as well.
Given $c\in X(n)$ we have that $f(c(\omega)) = f(c(\omega'))$ for all $\omega,\omega'\in\{0,1\}^n$
since $f\circ s^\ast = f\circ t^\ast$.
Thus we can just define $g_n(c) = f(c(0^n))$ which determines a map $g\colon X\to i_!S$ of cubesets with $g_0 = f$.
Uniqueness of $g$ is clear by construction.
\end{proof}

\begin{remark}
The set $\pi_0(X)$ can be computed as follows:
it is the set $X(0)$ modded out by the equivalence relation generated by
\[
    x\sim_0 y\iff\;\text{there exists}\; h\in X(1)\;\text{such that}\; s^\ast(h) = h(0) = x,\; t^\ast(h) = h(1) = y.
\]
If $X$ has the gluing property, then this relation already is an equivalence relation.

Moreover note that the coequalizer in the proposition is reflexive.
This gives another proof of the fact that the $0$-truncation, or $\pi_0$, commutes with finite products.
\end{remark}

\begin{definition}\label{def:ergodic}
Let $X$ be a cubeset.
We say that $X$ is ergodic if $\pi_0(X) = \ast$.
More generally, we say that $X$ is \emph{$n$-ergodic} if $\tau_{n-1}X = \ast$.
\end{definition}

\begin{remark}
If $X$ is concrete and fibrant, then $X$ is $n$-ergodic if and only if $X(n) = X(0)^{2^n}$.
Moreover, since the truncation preserves finite products by Proposition~\ref{prop:n-trunc-ihom-prod},
$n$-ergodic cubesets are closed under taking finite products.
\end{remark}

\begin{example}\label{ex:EM_k_ergodic}
The Eilenberg-Mac Lane cubeset $\cD_k(A)$ from Example~\ref{ex:host_kra_cubegroups} is $k$-ergodic
by Example~\ref{ex:truncation_HK}.
\end{example}

There also exists a decomposition into ergodic components.

\begin{proposition}\label{prop:ergodic-decomp}
Let $X$ be a cubeset and $\eta\colon X\to\tau_0X$ the unit of the $0$-truncation.
For $p\in\pi_0(X)$ denote by
\[
    X_p(n) = \{ c\in X(n) \mid \eta(c(\omega)) = p\;\text{for all}\;\omega\in\{0,1\}^n \}.
\]
Then $X_p$ is an ergodic sub-cubeset of $X$ and $X = \coprod_{p\in\pi_0(X)} X_p$.
\end{proposition}

\begin{proof}
It is clear that $X_p$ is a sub-cubeset of $X$ since for $c\in X_p(m)$ and $f\colon{\bbox}^n\to{\bbox}^m$
we have that $\eta(X(f)(c)(\omega)) = \eta(c(f(\omega))) = p$ for all $\omega\in\{0,1\}^n$
and thus $X(f)(c)\in X_p(n)$.
This induces a natural injective map $\coprod_{p\in\pi_0(X)} X_p\to X$.
It also is surjective: given $c\in X(n)$, to any two $\omega,\omega'\in\{0,1\}^n$
there is a map $l\colon{\bbox}^1\to{\bbox}^n$ with $l(0) = \omega$, $l(1) = \omega'$.
Then $c\circ l\in X(1)$ and by Proposition~\ref{prop:0-trunc-computation} we have that
\[
    \eta(c(\omega)) = \eta(c(l(0))) = \eta(c(l(1))) = \eta(c(\omega')).
\]
Lastly, $X_p$ is ergodic: given $x,y\in X_p(0)$, since $\eta(x) = \eta(y) = p$
we have that $x\sim_0 y$ in $X$.
This means that there is a finite sequence $h_1,\ldots,h_k\in X(1)$
such that $h_1(0) = x$, $h_i(1) = h_{i+1}(0)$ for $1\leq i\leq k-1$ and $h_k(1) = y$.
But then we have that $\eta(h_i(0)) = p = \eta(h_i(1))$ for all $i$,
and so $h_i\in X_p(1)$.
Hence $x\sim_0 y$ also in $X_p$ which means that $X_p$ is ergodic.
\end{proof}

We call every $X_p$ an \emph{ergodic component} of $X$
and the decomposition of $X$ into the coproduct $\coprod_{p\in\pi_0(X)}X_p$ the
\emph{decomposition in ergodic components}.

Let us also comment on connectivity of $\bbGamma$-sets.
For this we need the following propositions, characterizing connected $\bbGamma$-sets.
First, one can decompose $\bbGamma$-sets as a coproduct of the fibers over its points.

\begin{proposition}\label{prop:connected-comp-gamma}
Let $X$ be a $\bbGamma$-set, let $p\in X(0)$ and for $\langle n\rangle\in\bbGamma$ define
\[
    X_p(n) = \{ c\in X(n) \mid c(0) = p \}.
\]
Then $X_p$ is a sub-$\bbGamma$-set of $X$ and $X = \coprod_{p\in X(0)} X_p$.
This is the decomposition of $X$ into its connected components.
\end{proposition}

\begin{proof}
Since there is a unique map $u\colon\ast\to\langle n\rangle$, $X_p$ is a $\bbGamma$-set, and it is clear that $X_p\sub X$.
It is clear that $X_p$ as a $\bbGamma$-set is connected since every $\bbGamma$-set with $X(0) = \ast$ is connected.
Lastly, $X(n) = \coprod_{p\in X(0)} X_p(n)$ is clear by definition, and so $X = \coprod_{p\in X(0)} X_p$.
\end{proof}

\begin{definition}
Let $\cC$ be a category.
An object $C\in\cC$ is called \emph{connected} if $\hom(C,-)\colon\cC\to\set$ preserves coproducts.
\end{definition}

\begin{proposition}\label{prop:connected-gamma-prop}
Let $X$ be a $\bbGamma$-set.
Then the following conditions are equivalent.
\begin{enumerate}[(a)]
    \item $X$ is connected.
    \item For every decomposition $X = X_1\sqcup X_2$ we have that $X_1$ is empty or equal to $X$.
    \item $X(0)$ is a singleton.
\end{enumerate}
\end{proposition}

\begin{proof}
$(a)\implies(b)$: Let $X = X_1\sqcup X_2$.
Then the identity on $X$ has to factor w.l.o.g. as $X\to X_1\to X$.
Therefore, $X_1\to X$ is an epimorphism, and hence an isomorphism.

$(b)\implies(c)$: Follows from Proposition~\ref{prop:connected-comp-gamma}.

$(c)\implies(a)$: If $X(0)$ is a point, then for any map $X\to\coprod_i Y_i$
its underlying point factors through one $Y_i$,
and so does every map $X(n)\to\coprod_i Y_i(n)$.
\end{proof}

Therefore, we have a decomposition of the $\bbGamma$-set $X$ into its connected components:
it is a decomposition $X = \coprod_{i\in I} X_i$ where $X_i$ are connected and $X_i\sub X$ are sub-$\bbGamma$-sets.

\subsubsection{Truncations over \texorpdfstring{$\bbGamma$}{Gamma}}\label{ssec:truncation-gamma}

Note that of course one can also define truncations over $\bbGamma$ instead of $\bbox$.
Here the correct set of morphisms to localize at is the refined $k$-corners for $k>n$ from Section~\ref{ssec:corner-gamma}.
Recall the set
\[
    \Lambda_{\bbGamma} \coloneq \{{\bbcor}^n_\omega\hookrightarrow{\bblox}^n \mid n\in\N,\,\omega\in\{0,1\}^n\setminus\{0^n\}\}
\]
from Definition~\ref{def:corner_in_gamma} and~\ref{def:omega_corner_gamma}.
This yields the variant
\[
    \Lambda_{\bbGamma,n} \coloneq \{{\bbcor}^k_\omega\hookrightarrow{\bblox}^k \mid k\geq n+1,\,\omega\in\{0,1\}^k\setminus\{0^k\}\}.
\]
of maps in $\Lambda_{\bbGamma}$ between objects of dimension $>n$.
Then we can define the truncations analogously to $\cSet$.

\begin{definition}
A $\bbGamma$-set $X$ is called $n$-step if it is local w.r.t. $\Lambda_{\bbGamma,n}$.
We denote by $\PSh(\bbGamma)_n$ the full subcategory of $n$-step $\bbGamma$-sets
and by $\tau_n^{\bbGamma}\colon\PSh(\bbGamma)\to\PSh(\bbGamma)_n$ the localization onto $\Lambda_{\bbGamma,n}$-local objects.
\end{definition}

\begin{remark}
Essentially all the same properties of the truncation that work with cubesets hold over $\bbGamma$ as well,
e.g., the subcategory of $n$-step $\bbGamma$-sets is closed under limits, filtered colimits and coproducts.
\end{remark}

Again, one can compute the truncation similarly to cubesets.
However, note that the truncation does not affect the value $X(0)$ at all:
the $\bbGamma$-set decomposes into a coproduct of its $X(0)$-values, see Proposition~\ref{prop:connected-comp-gamma},
and so it suffices to see that $\tau_n X(0) = \ast$ whenever $X(0) = \ast$.
But this follows from the explicit construction of the reflection $\tau_n X$ as a filtered colimit over objects which arise as connected colimit of objects $Y$ that all have $Y(0) = \ast$,
see \cite[Construction 1.37]{Adamek1994}.

But other than that, for a $\bbGamma$-set arising from a concrete fibrant cubeset, computing the truncation over $\bbGamma$ is analogous to computing it over $\bbox$.
This will need the following Definition.

\begin{definition}
Let $X$ be a concrete fibrant cubeset.
For $n=0$ we define $x\approx_n y$ for $x,y\in X(1)$ if $x(0) = y(0)$.
For $n>0$ we define $x\approx_n y$ for $x,y\in X(1)$ if there exist $\omega\in\{0,1\}^{n+1}\setminus\{0^{n+1}\}$ and $c,c'\in X(n+1)$ such that
\[
    c\rvert_{j_!{\bbcor}^{n+1}_\omega} = c'\rvert_{j_!{\bbcor}^{n+1}_\omega} \quad\text{and}\quad c(0,\omega) = x,\; c'(0,\omega) = y.
\]
\end{definition}

Here $c(0,\omega)$ denotes the composition $c\circ t$ where $t\colon{\bblox}^1\to{\bblox}^{n+1}$ is the unique linear map with $t(1) = \omega$.

\begin{proposition}
Let $X$ be a concrete fibrant cubeset, $n\in\N$ and $x,y\in X(1)$.
Then the following assertions are equivalent.
\begin{enumerate}[(a)]
    \item $x\approx_n y$.
    \item $x(0) = y(0)$ and $x(1)\sim_n y(1)$.
\end{enumerate}
In particular, the relation $\approx_n$ is an equivalence relation on $X(1)$ for all $n\in\N_0$.
\end{proposition}

\begin{proof}
$(a)\implies (b)$:
If $x\approx_n y$, there exist $\omega\in\{0,1\}^{n+1}\setminus\{0^{n+1}\}$ and $c,c'\in X(n+1)$ such that
\[
    c\rvert_{j_!{\bbcor}^{n+1}_\omega} = c'\rvert_{j_!{\bbcor}^{n+1}_\omega} \quad\text{and}\quad c(0,\omega) = x,\; c'(0,\omega) = y.
\]
In particular, $x(0) = y(0)$.
Moreover, we have seen in Lemma~\ref{lem:alternative-nilfibrant-concrete}
that $X$ is local w.r.t. $j_!{\bbcor}^k_\omega\to{\bbcor}^k_\omega$.
This means that $c\rvert_{{\bbcor}^{n+1}_\omega} = c'\rvert_{{\bbcor}^{n+1}_\omega}$.
Now take an automorphism of cubes exchanging $\omega$ and $1^{n+1}$ (using the appropriate flips).
This implies that $c\rvert_{{\bbcor}^{n+1}_{1^{n+1}}} = c'\rvert_{{\bbcor}^{n+1}_{1^{n+1}}}$ and in particular since ${\bbcor}^{n+1}\sub{\bbcor}^{n+1}_{1^{n+1}}$,
$c\rvert_{{\bbcor}^{n+1}} = c'\rvert_{{\bbcor}^{n+1}}$ and $c(1^{n+1}) = x(1)$, $c'(1^{n+1}) = y(1)$.
This means that $x(1)\sim_n y(1)$.

$(b)\implies (a)$:
If $x(1)\sim_n y(1)$ take the (linear!) degeneracy $p\colon{\bbox}^{n+1}\to{\bbox}^1$ along the $x_1$-coordinate axis.
Then $c = x\circ p$ is an ($n+1$)-cube in $X$ with $c(1^{n+1}) = x(1)$.
Using universal replacement, we see that the configuration
\[
    c'\colon\{0,1\}^{n+1}\to X(0),\quad\omega\mapsto\begin{cases}
        c(\omega),\quad &\omega\neq 1^{n+1}, \\
        y(1),\quad &\omega = 1^{n+1},
    \end{cases}
\]
is in $X(n+1)$.
Thus we have that
\[
    c\rvert_{j_!{\bbcor}^{n+1}_{1^{n+1}}} = c'\rvert_{j_!{\bbcor}^{n+1}_{1^{n+1}}} \quad\text{and}\quad c(0,\omega) = x,\; c'(0,\omega) = y.\qedhere
\]
\end{proof}

This equivalence relation yields a $\bbGamma$-set $X/\approx_n$ as follows:
for $m=0$ we set $X/\approx_n(0) = X(0)$ and for $m>0$ we define $X/\approx_n(m)$ to be the quotient of $X(m)$ that identifies two $m$-cubes $a,b$
if for all inclusions $l\colon{\bblox}^1\to{\bblox}^m$ we have that $a\circ l \approx_n b\circ l$.

The previous proposition gives us an easy description of $X/\approx_n$:
two $m$-cubes $a$ and $b$ are identified if and only if $a(0) = b(0)$ and $a(\omega) \sim_n b(\omega)$ for all $\omega\in\{0,1\}^m\setminus\{0^m\}$.
Our goal now is to show that this quotient $X/\approx_n$ actually is the truncation over $\bbGamma$.

\begin{proposition}\label{prop:trunc-concfib-gamma}
Let $X$ be a concrete fibrant cubeset and $n\in\N_0$.
Then the $\bbGamma$-truncation of degree $n$ of $j^\ast X$ is $X/\approx_n$ and the unit $j^\ast X\to\tau^{\bbGamma}_n j^\ast X$ is a $\bbGamma$-fibration.
\end{proposition}

\begin{proof}
The proof works essentially the same as the proof of Proposition~\ref{prop:trunc_concfib}.
In a first step, we show existence of a map making the diagram
\begin{center}
\begin{tikzcd}
	{j^\ast X} && {\tau^{\bbGamma}_n j^\ast X} \\
	& {X/\approx_n}
	\arrow["\eta", from=1-1, to=1-3]
	\arrow["\pi"', from=1-1, to=2-2]
	\arrow["f"', dashed, from=2-2, to=1-3]
\end{tikzcd}
\end{center}
commutative.
We compute $f$ sectionwise, and $k = 0$ is clear since $\tau^{\bbGamma}_n j^\ast X(0) = X(0) = X/\approx_n(0)$.
For $1\leq k\leq n+1$ we need to see that if $a,b$ are two $k$-cubes in $X$ with $a\approx_n b$, then $\eta(a) = \eta(b)$.
Take a linear surjection (a degeneracy) $p\colon{\bbox}^{n+1}\to{\bbox}^k$.
Then $a' = p^\ast(a) \approx_n p^\ast(b) = b'$.
Using universal replacement we can, one by one, exchange each point of $a'$, which is not at position $0$, by a point of $b'$ at the same position.
But instead of rotating $a'$, we now use unique corner completion of $\tau_n^{\bbGamma} j^\ast X$ against every $\omega$-corner ${\bbcor}^{n+1}_\omega$
to conclude that $\eta(a') = \eta(c_k)$ for each $c_k$ after a finite replacement, and finally, $\eta(a) = \eta(b)$.
For $k\geq n+2$ the same induction as in the proof of Proposition~\ref{prop:trunc_concfib} yields the claim.

The rest of the proof works verbatim as in the proof of Proposition~\ref{prop:trunc_concfib}
except that at $\omega = 0$, equality has to occur in the equivalence relation, and one has to consider all corners ${\bbcor}^k_\omega$.

For $n=0$ note that the Hom-isomorphism of the adjunction holds since both sides agree with $Y(0)^{X(0)}$.
\end{proof}

